\section*{Chapter \ref{ch:b3:mantle-convection} --- Mantle Convection and Plate Driving Forces}

\begin{solution}{exo:b3:mantle-convection:1}
$\mathrm{Ra} = 3500\times9.9\times\num{3e-5}\times1000\times(\num{6.7e5})^3/
(10^{-6}\times\num{5e20}) = \num{6.25e5}$, about 950 times the critical
657.5: the layer convects vigorously.
\end{solution}

\begin{solution}{exo:b3:mantle-convection:2}
$\tau_M = 10^{21}/\num{7.7e10} = \qty{1.30e10}{s}$, about \qty{412}{yr}. The
tide, which lasts hours, meets an elastic mantle; the melting of an ice sheet,
over thousands of years, and the sinking of a slab, over millions, meet a
viscous one.
\end{solution}

\begin{solution}{exo:b3:mantle-convection:3}
$\tau = 4\pi\times10^{21}/(3300\times9.8\lambda)$: \qty{3.89e11}{s} =
\qty{12.3}{kyr} for \qty{1000}{km}, and \qty{4.1}{kyr} for \qty{3000}{km}. The
broader load recovers faster.
\end{solution}

\begin{solution}{exo:b3:mantle-convection:4}
$g\rho_m\alpha(T_m - T_0)\kappa t = 9.8\times3300\times\num{3.1e-5}\times1450
\times\num{0.8e-6}\times50\times\num{3.156e13} = \qty{1.835e12}{N.m^{-1}}$; the
bracket is $1 + 2\times3300\times\num{3.1e-5}\times1450/(\pi\times2270) = 1.0416$,
so $F_{rp} = \qty{1.91e12}{N.m^{-1}}$, 24 times smaller than the slab pull.
\end{solution}

\begin{solution}{exo:b3:mantle-convection:5}
(a) $v = 3300\times9.8\times\num{3e-5}\times300\times(\num{e4})^2/10^{21} =
\qty{2.91e-11}{m.s^{-1}}$; the rise time across \qty{10}{km} is $d/v =
\qty{3.44e14}{s}$, \qty{10.9}{Myr}, the diffusion time $d^2/\kappa =
\qty{e14}{s}$, \qty{3.2}{Myr}. Their ratio is $(d^2/\kappa)/(d/v) =
vd/\kappa = \rho g\alpha\Delta Td^3/(\kappa\eta) = \mathrm{Ra} = 0.29$. (b) It
cools first: small blobs do not convect. (c) $\mathrm{Ra} \propto d^3$, so $d =
10\times(657.5/0.291)^{1/3} = \qty{131}{km}$.
\end{solution}

\begin{solution}{exo:b3:mantle-convection:6}
$2\mathrm{Ra}_c = 1315$ and $\mathrm{Ra}/2\mathrm{Ra}_c = 7604$, whose cube
root is 19.66. $\delta = 2891/19.66 = \qty{147}{km}$, $\mathrm{Nu} = 19.66/2 =
9.8$, $u = (10^{-6}/\num{2.891e6})\times19.66^2 = \qty{1.34e-10}{m.s^{-1}} =
\qty{0.42}{cm/yr}$, and $q = 9.8\times3.1\times1450/\num{2.891e6} =
\qty{15}{mW.m^{-2}}$.
\end{solution}

\begin{solution}{exo:b3:mantle-convection:7}
$u = 0.09/\num{3.156e7} = \qty{2.85e-9}{m.s^{-1}}$ and $B = 2270\times\num{8e8}
\times\num{2.85e-9} = \qty{5.18e3}{kg.s^{-1}} = \qty{5.2}{Mg/s}$, about
three-fifths of Hawaii's. The heat flux is $1170\times5178/\num{3.1e-5} =
\qty{1.95e11}{W}$, \qty{0.2}{TW}, or \qty{0.42}{\percent} of the Earth's heat
loss.
\end{solution}

\begin{solution}{exo:b3:mantle-convection:8}
$v = 2\times30\times9.8\times(\num{4e5})^2/(9\times10^{21}) =
\qty{1.05e-8}{m.s^{-1}} = \qty{0.33}{m/yr}$; \qty{2000}{km} take
$\num{2e6}/\num{1.05e-8} = \qty{1.9e14}{s}$, about \qty{6}{Myr}.
\end{solution}

\begin{solution}{exo:b3:mantle-convection:9}
$w = 300e^{-10/4.4} = \qty{30.9}{m}$ and $\dot w = w/\tau = 30.9/4400 =
\qty{7.0}{mm/yr}$, two-thirds of the rate measured near Umeå: a uniform
half-space with this $\tau$ is in the right range, but the real rebound is
faster or its depression deeper.
\end{solution}

\begin{solution}{exo:b3:mantle-convection:10}
(a) With $w = 2\rho_m\alpha\Delta T\sqrt{\kappa t/\pi}/(\rho_m - \rho_w)$,
\[
  \tfrac12(\rho_m - \rho_w)gw^2 = \frac{2g\rho_m^2\alpha^2\Delta T^2\kappa t}
  {\pi(\rho_m - \rho_w)} = g\rho_m\alpha\Delta T\kappa t\times
  \frac{2\rho_m\alpha\Delta T}{\pi(\rho_m - \rho_w)} .
\] (b) By parts,
$\int_0^\infty x\operatorname{erfc}x\,\mathrm dx = [\tfrac{x^2}2\operatorname{erfc}
x]_0^\infty + \int_0^\infty\tfrac{x^2}2\,\tfrac2{\sqrt\pi}e^{-x^2}\,\mathrm dx =
0 + \tfrac1{\sqrt\pi}\cdot\tfrac{\sqrt\pi}4 = \tfrac14$. (c) At \qty{100}{Myr},
$\rho_mg\alpha\int z\theta = \qty{3.67e12}{N.m^{-1}}$ and, with $w =
\qty{3705}{m}$, $\tfrac12(\rho_m - \rho_w)gw^2 = \qty{1.53e11}{N.m^{-1}}$: the
water term is only \qty{4}{\percent} of the force.
\end{solution}

\begin{solution}{exo:b3:mantle-convection:11}
$\ln f = 3\ln(\pi^2 + k^2) - 2\ln k$, whose derivative $6k/(\pi^2 + k^2) - 2/k$
vanishes when $3k^2 = \pi^2 + k^2$, $k^2 = \pi^2/2$. There $f = (3\pi^2/2)^3/
(\pi^2/2) = 27\pi^4/4$. The wavelength $2\pi/k = 2\sqrt2\,d = 2.83\,d$ holds a
pair of counter-rotating cells, each $\sqrt2\,d = 1.41\,d$ wide.
\end{solution}

\begin{solution}{exo:b3:mantle-convection:12}
(a) $u = \qty{1.90e-9}{m.s^{-1}}$, $\eta u/h = 10^{20}\times\num{1.90e-9}/
\num{1.5e5} = \qty{1.27}{MPa}$, and over \qty{6000}{km}, \qty{7.6e12}{N.m^{-1}}.
(b) Twice the ridge push and a sixth of the slab pull: the drag can resist
the ridge push alone, but not the slab. (c) $\eta uL/h = \num{3.8e12} +
\num{4.6e12}$ gives $u = \num{8.4e12}\times\num{1.5e5}/(\num{6e6}\times10^{20})
= \qty{2.1e-9}{m.s^{-1}} = \qty{6.6}{cm/yr}$, the speed of the fast oceanic
plates.
\end{solution}

\begin{solution}{pb:b3:mantle-convection:1}
\textbf{1.} $p = 1000\times9.8\times300 = \qty{2.94}{MPa}$.

\textbf{2.} $w = \rho_wh/\rho_m = 300\times1000/3300 = \qty{90.9}{m}$. The crust
moves down as a block; what flows out from beneath the load, and back, is the
mantle, so the column is balanced by mantle displaced at depth.

\textbf{3.} $\pi\times(\num{1.25e5})^2 = \qty{4.91e10}{m^2}$, and the water
weighs $\num{4.91e10}\times300\times1000 = \qty{1.47e16}{kg}$.

\textbf{4.} The shoreline formed level, when the floor was depressed most at
the centre; the centre has risen more than the rim since, so the line now
climbs towards the centre, by \qty{77.2}{m} over \qty{125}{km}:
\qty{0.62}{m/km}.

\textbf{5.} $\lambda/2 = \qty{250}{km}$, so $\lambda = \qty{500}{km}$ and $k =
2\pi/\num{5e5} = \qty{1.257e-5}{m^{-1}}$.

\textbf{6.} $1/k = \lambda/2\pi = \qty{79.6}{km}$.

\textbf{7.} $w = 90.9 - \text{uplift}$: 79.4, 69.3, 52.9, 40.4, 26.9 and
\qty{13.7}{m}.

\textbf{8.} $w = w_0e^{-t/\tau}$ gives $\ln w = \ln w_0 - t/\tau$, a straight line
of slope $-1/\tau$.

\textbf{9.} $\tau = (14 - 1)/\ln(79.4/13.7) = 13/1.7571 = \qty{7.40}{kyr}$.

\textbf{10.} $w(6) = 79.4\,e^{-5/7.40} = \qty{40.4}{m}$, an uplift of
\qty{50.5}{m}, as in the table.

\textbf{11.} \qty{13.7}{m} are left; $\dot w = 13.7/7400 = \qty{1.85}{mm/yr}$.

\textbf{12.} The whole \qty{13.7}{m}, approached ever more slowly.

\textbf{13.} The rate falls as $e^{-t/\tau}$: $t = 7.40\ln(1.85/0.5) =
7.40\times1.309 = \qty{9.7}{kyr}$, some 9700 more years.

\textbf{14.} $\tau = 7399\times\num{3.156e7} = \qty{2.335e11}{s}$, and $\eta =
\tau\rho_mg\lambda/(4\pi) = \num{2.335e11}\times3300\times9.8\times\num{5e5}/
12.566 = \qty{3.0e20}{Pa.s}$.

\textbf{15.} Just below the lower end of the Fennoscandian range. The flow
beneath a load \qty{500}{km} in wavelength decays over \qty{80}{km}: the
rebound samples the shallowest, weakest mantle, beneath and partly within what
is, in the real Earth, the lithosphere.

\textbf{16.} $\tau = 4\pi\times\num{3.0e20}/(3300\times9.8\times\num{3e6}) =
\qty{3.9e10}{s} = \qty{1.2}{kyr}$. Beneath a load that broad the flow reaches
down \qty{477}{km} and more, towards a lower mantle ten times more viscous,
which slows it.

\textbf{17.} It holds part of a narrow load by flexure: the depression at
equilibrium is shallower and wider than the isostatic \qty{90.9}{m}, and part of
the rebound is elastic and immediate. A uniform viscous half-space fitted to
the uplift then gives a viscosity that mixes both effects.

\textbf{18.} $\tau_M = \num{3.0e20}/\num{7.7e10} = \qty{3.9e9}{s}$, about
\qty{124}{yr}, sixty times shorter than $\tau$: over thousands of years the
mantle responds as a viscous fluid.

\textbf{19.} $1 - e^{-50/7.40} = 0.9988$: \qty{99.9}{\percent} of the
equilibrium depression was reached, and the assumption holds.

\textbf{20.} $\eta \propto \lambda$: halving $\lambda$ halves $\eta$, and
$\log_{10}\eta$ falls by $\log_{10}2 = 0.30$.

\textbf{21.} $\mathrm{Ra} = 3700\times9.9\times\num{3.1e-5}\times1450\times
(\num{6.7e5})^3/(10^{-6}\times\num{3.0e20}) = \num{1.65e6}$.

\textbf{22.} $\num{1.65e6}/657.5 \approx 2500$ times critical.

\textbf{23.} $(\mathrm{Ra}/2\mathrm{Ra}_c)^{1/3} = (1254)^{1/3} = 10.78$: $\delta
= 670/10.78 = \qty{62}{km}$ and $\mathrm{Nu} = 5.4$.

\textbf{24.} $u = (10^{-6}/\num{6.7e5})\times10.78^2 = \qty{1.73e-10}{m.s^{-1}}
= \qty{0.55}{cm/yr}$, two to fourteen times slower than the plates.

\textbf{25.} $q = 5.4\times3.1\times1450/\num{6.7e5} = \qty{36}{mW.m^{-2}}$. The
layer leaves out the heat produced within the mantle and the heat from the
core and the lower mantle; and the scaling drops its constants of order one,
while the plates, cooling as half-spaces from the ridges, lose far more heat on
young floor than a mean boundary layer would.

\textbf{26.} $\eta = \qty{3.0e20}{Pa.s}$: $\log_{10}\eta =
\textbf{\num[text-series-to-math=true]{20.48}}$.
\end{solution}
